%% highlights.tex — CVIU 强制要件：独立可编辑文件（文件名须含 "highlights"）
%% 规格: 3–5 条, 每条 ≤85 字符含空格 — 五条已程序化实测 65–85 字符（reports/submission-package-2026-08-26.md §2.2）
%% 维护纪律: 与 main.tex 版内 highlights 环境同步改；数字若因预注册条款回改须联动刷新。

\documentclass{article}
\usepackage[margin=1in]{geometry}
\pagestyle{empty}
\begin{document}
\section*{Highlights}
%% 2026-11-15 深度改造: 与 main.tex 版内 highlights 同步（R2-1/R2-5 联动）
%% 2026-11-16 统稿: 第4条同步 main.tex 的头对头基线句（按 ≤85 字符约束压缩表述, 实测 83 字符; 主文为完整版）。
\begin{itemize}
\item Decoupled framework absorbs taxonomy evolution cheaply (synthetic tier)
\item Taxonomy transitions cost $\geq$3x less wall clock (measured 6.07x, synthetic tier)
\item Warm start: 82.0\% top-1 from 20 clips (vs.\ 96.6\% full-budget, offset tier)
\item NTU60 frozen features retain 90.7\% at 10\% labels, above AimCLR and ST-GCN baselines
\item Canine tier stays below chance at 13\% budget (9.8\% vs.\ 11.1\%)---honest boundary
\end{itemize}
\end{document}
